% id: 51-26
% book: 쎈B 공통수학2 · 04 도형의 이동
% section: 유형 05 점의 대칭이동
% figure: none
% answer: 
% answer_source: 
% variant_level: 0 (원본 전사)
% 이 파일은 build.mjs 가 items.json 에서 만든다. 고칠 때는 items.json 을 고친다.
\smash{\raisebox{1.5mm}{\dmkichul{쎈B 공통수학2 51쪽 26번}}}
점~$\pt{P}_1(3{,}\mkern3mu\mathopen{}1)$을 원점에 대하여 대칭이동한 점을 $\pt{P}_2$라 하고, 점~$\pt{P}_2$를 직선 $y=x$에 대하여 대칭이동한 점을 $\pt{P}_3$이라 하자. 다시 점~$\pt{P}_3$을 원점에 대하여 대칭이동한 점을 $\pt{P}_4$라 하고, 점~$\pt{P}_4$를 직선 $y=x$에 대하여 대칭이동한 점을 $\pt{P}_5$라 하자. 이와 같은 방법으로 원점과 직선 $y=x$에 대하여 대칭이동한 점을 차례대로 $\pt{P}_6$, $\pt{P}_7$, $\cdots$이라 할 때, 점~$\pt{P}_{111}$의 좌표를 구하시오.
